\documentclass[11pt]{article}

\usepackage[margin=1.1in]{geometry}
\usepackage[utf8]{inputenc}
\usepackage[T1]{fontenc}
\DeclareUnicodeCharacter{03B8}{\ensuremath{\theta}}
\DeclareUnicodeCharacter{2212}{\ensuremath{-}}
\usepackage{amsmath,amssymb}
\usepackage{booktabs}
\usepackage{microtype}
\usepackage[hidelinks]{hyperref}

\newcommand{\E}{\mathbb{E}}
\newcommand{\code}[1]{\texttt{#1}}

\title{Concourse: Four Design Questions\\
\large Background and questions for the next design conversation}
\author{Concourse research notes}
\date{July 17, 2026}

\begin{document}
\maketitle
\tableofcontents

%==============================================================================
\section{What this document is}

Concourse now simulates every queueing model we set out to support, and it
computes parameter derivatives for most of them. Three capabilities remain
open, and behind them sits one deeper question about how the system records
its history. This document explains the background needed to think about
those questions, then states each question in plain terms.

Each section builds on the ones before it. The background sections describe
how the simulator works, what it writes down, and how derivative estimation
uses what it wrote down. The question sections then say what is missing,
why it is missing, and what a decision would need to settle.

%==============================================================================
\section{Background}

\subsection{How the simulator runs a queueing model}

A queueing model describes jobs that arrive, wait, get served, and leave.
The simulator does not think in terms of jobs directly. It thinks in terms
of \emph{clocks}.

A clock is a pending event. When a job is being served, there is a clock
for the moment its service will finish. When a source is active, there is a
clock for the next arrival. Each clock has two pieces of information: a
probability distribution for how long it will run, and the time it was
started. We call the start time the enabling time.

The simulation is a race. All the active clocks run at once. The one that
fires first wins. Firing a clock changes the state: a job moves, a queue
grows or shrinks, a job leaves the system. The change can also start new
clocks and cancel old ones. Then the race resumes among the clocks that
remain. This loop is the entire simulation.

One rule governs everything the model tells the sampler. The model's whole
statement about a clock is the pair (distribution, enabling time). The
model never tells the sampler \emph{how} to draw random numbers. It only
declares what the law of the firing time is, given the history so far. The
sampler is free to implement that law any way it likes.

\subsection{The parameter vector}

Models have tunable numbers: an arrival rate, a service rate, a patience
rate. We collect them in a vector called $\theta$. The value of $\theta$
enters the simulation in exactly one place: when the model builds a clock's
distribution. Nothing else in the system reads $\theta$. This discipline is
what makes derivative estimation possible later, because there is only one
door through which $\theta$ can affect anything.

\subsection{The record}

Each run writes down a record. The record is a list with one entry per
firing. An entry holds three things: which clock fired, the time it fired,
and any extra random values that the firing consumed (more on those below).

The record is the primary output of a run. Every statistic we report is
computed by replaying the record: start from the initial state, apply the
firings one at a time, and watch the state evolve. Replaying is cheap and
it is exact, because the state change caused by a firing is a deterministic
function once the recorded random values are given.

\subsection{Auxiliary draws}

Most randomness lives in the clocks. Two kinds do not.

First, \emph{marks}. When a job arrives it may get random attributes: a
size, a class. These are drawn at the moment of arrival, not raced as
clocks.

Second, \emph{random routing}. Some models flip a weighted coin to decide
where a finished job goes next.

We call these auxiliary draws. The rule for them is strict: every
auxiliary draw is taken from a named random stream and written into the
record. Replay never draws again; it reads the recorded values back. This
is what keeps replay exact even for models with marks and coin flips.

\subsection{Re-declared clocks and segments}

Usually a clock keeps one distribution for its whole life. Processor
sharing breaks this pattern, and it is worth understanding why, because two
of the open questions trace back to it.

Under processor sharing, all the jobs at a station are served at the same
time, each at a fraction of the server's speed. When a new job arrives, the
fraction shrinks for everyone. When a job finishes, the fraction grows.
So the law of each remaining service time changes every time the
number of jobs changes, even though the clocks themselves stay active.

We handle this by letting the model \emph{re-declare} a clock. The clock
keeps its original enabling time. What changes is the declared
distribution, which now says: given how much work this job has already
received, and given the current sharing fraction, here is the law of the
remaining time. Each re-declaration starts a new \emph{segment} in the
clock's history. A clock that was never re-declared has one segment. A
processor-sharing clock can have many.

Segments matter because the derivative machinery replays clock histories,
and a multi-segment history is genuinely more complicated than a
single-segment one.

\subsection{Why we want derivatives}

Suppose the average waiting time is $W(\theta)$. We often want
$\partial W / \partial \theta$: how much the waiting time would change if
we made service ten percent faster. Derivatives drive optimization,
sensitivity analysis, and calibration against data.

The difficulty is that $W$ is estimated by simulation. Running the model
twice at two nearby values of $\theta$ and subtracting does not work well:
the noise of the two runs swamps the small difference between them, unless
the runs share their randomness in a carefully controlled way. The field has developed several estimators that compute
the derivative from simulation output directly. Concourse supports four,
through the ClockGradients package. They differ in what they need and in
what they can handle.

\subsection{The four estimators, in one paragraph each}

\emph{Score function.} Replay the record and compute how much more or less
likely the same trajectory would have been at a slightly different
$\theta$. Multiply that sensitivity by the observed outcome and average
over runs. This works for almost any model and any statistic. Its weakness
is noise: the estimates have high variance.

\emph{Pathwise, also called IPA.} Keep the order of firings fixed, and ask
how each firing \emph{time} would shift if $\theta$ moved a little. Slide
the times, recompute the statistic, and read off the slope. This has low
noise. Its weakness is blindness: if changing $\theta$ would change which
clock wins a race, IPA does not see that effect, and for statistics that
jump when the order changes, IPA is simply wrong.

\emph{Branching.} Run the simulation live. At a firing, make a copy of the
whole running world and force a \emph{different} clock to win in the copy.
Compare the outcome of the copy with the outcome of the original. This
directly measures the effect of order changes, which is exactly what IPA
misses. Its weakness is cost: every branch point spawns clones.

\emph{Smoothed perturbation analysis, SPA.} A compromise. Take the IPA
estimate, then add a correction for the near-ties: the moments where two
clocks almost fired at the same time, so a small change in $\theta$ could
swap their order. The correction weighs each near-tie by how likely the
swap was, using the hazard of the runner-up clock, and by how much the
statistic would jump if the swap happened. It needs fewer clones than
branching.

\subsection{Cloning a world}

Branching and SPA need to copy a running simulation. The copy must be
exact: the same state, the same clocks, the same ages, and the same
pending randomness. If you let the original and the copy both run forward
with no other action, they must produce identical futures, firing for
firing. We call this a coupled clone.

Sometimes we want the opposite. After cloning, we give the copy fresh
randomness derived from a chosen seed. Then the copy's future is a fresh
draw, independent of the original. And if two copies are given the
\emph{same} seed, their futures match each other while differing from the
original. This controlled sharing of randomness is what makes the
difference between two clones quiet: the shared randomness cancels, and
only the effect we care about remains.

Concourse implements all of this today, with one restriction that leads
directly to the first question.

%==============================================================================
\section{Question 1: cloning worlds that use auxiliary draws}

\subsection{Where we are}

The branchable world currently refuses any model that uses marks or random
routing. The refusal is deliberate. Replay of a record is safe for these
models, because the recorded values are read back. But a clone runs
\emph{forward}, past the end of any record. Its future arrivals need new
sizes. Its future routing needs new coin flips. There is nothing recorded
to read back, so the world would need live random streams for auxiliary
draws, and we have not yet decided how those streams should behave under
cloning.

\subsection{What has to be decided}

The clock randomness already has good cloning behavior, because the
sampler keys each clock's stream by the clock's name. The auxiliary
streams need the same three properties, and each one is a design choice.

First, \emph{copying}. When a world is cloned, the auxiliary streams must
be copied with their current positions, so that the clone's next size draw
equals the original's next size draw. Otherwise the coupled clone would
drift at the first arrival.

Second, \emph{re-seeding}. When a clone is given fresh randomness, the
auxiliary streams must be re-seeded too. If we re-seed only the clocks,
the clone would replay the original's future sizes, and the two would stay
partly correlated when they should be independent.

Third, and hardest, \emph{alignment}. Two same-seed clones are supposed to
see matching randomness. But after a forced firing, the two clones can
disagree about the order of later events. Suppose clone A sees its next
two arrivals in the order job then job, while clone B, whose forced firing
delayed something, sees an extra service completion in between. If size
draws are handed out in the order they are requested, the two clones may
consume the stream at different points, and their jobs stop matching. The
design choice is what the draws should be keyed by. Keying by request
order is simple and wrong. Keying by something stable, such as the
arrival's position in that source's own sequence, keeps the fifth arrival
of clone A matched with the fifth arrival of clone B, no matter what else
happened in between.

There is also a smaller, separable question. If a mark's distribution
depends on $\theta$, the derivative of the mark itself becomes part of the
answer, and the estimators need an extra term for it. Everything is
simpler when marks do not depend on $\theta$. We can choose to support
$\theta$-free marks first, and treat $\theta$-dependent marks as a later
step.

\subsection{What a decision looks like}

A decision here settles: where the auxiliary streams live (in the world
object, alongside the sampler), what re-seed covers, and what the stream
key is for each kind of draw. With those three fixed, the implementation
is mechanical, and the certification harness can be extended to test the
new obligations: clone-copies-marks, reseed-covers-marks, and same-seed
clones agree on the sizes of corresponding jobs.

%==============================================================================
\section{Question 2: SPA's correction for re-declared clocks}

\subsection{Where we are}

SPA works today for models whose clocks keep one distribution for life. It
refuses, with a clear message, any record that contains a multi-segment
clock. So SPA runs on our FCFS and priority models, but not on processor
sharing. The refusal is not an implementation gap. The mathematics of the
correction term has not been worked out for the multi-segment case.

\subsection{The problem in plain terms}

Recall what the correction does. At a near-tie, two clocks could have
fired in either order. The correction estimates two things: the
probability rate at which the order would swap as $\theta$ moves, and the
change in the statistic if it did swap.

The swap rate is built from the runner-up clock's hazard rate at the
decision time. The hazard rate answers: given that this clock has survived
until now, how likely is it to fire in the next instant. For a
single-segment clock, the hazard comes straight from the one declared
distribution and the clock's age.

For a re-declared clock, the current law is the latest segment's law, and
the age must be measured against the right anchor. That part is not deep;
we know the current declared law at every moment. The deeper issue is the
paired-clone construction around the swap. The two clones continue the
world after firing the pair in both orders, and the correctness argument
for the correction assumes the surviving clock's remaining life is
conditioned in a specific way. When the surviving clock can be re-declared
by the very firing we forced --- and under processor sharing it always is,
because forcing a service completion changes the sharing fraction for
everyone --- the existing argument no longer covers the construction. The
weight, the conditioning, and possibly the variance of the estimator all
need to be re-derived.

\subsection{What has to be decided}

Two things, one mathematical and one structural.

The mathematical piece: derive the boundary weight when the runner-up has
a multi-segment history, and check whether the paired-clone construction
still cancels the shared prefix, or needs modification. This is research,
and a small hand-worked example (two jobs in a processor-sharing station)
would be the right first step. It could confirm the natural guess, or
show that a new term appears.

The structural piece: the refusal currently happens for the whole record.
A record could contain many single-segment candidates and only a few
re-declared ones. Once the mathematics is settled we should decide whether
validity is a property of the whole record or of each candidate
separately, and place the check where the decision says it belongs.

%==============================================================================
\section{Question 3: IPA over processor-sharing records}

\subsection{Where we are}

The IPA machinery already supports multi-segment records. The replay
formulas that slide firing times around were written with segments in
mind, and the score estimator, which shares the same record structure,
passes its tests on processor sharing. What we have not done is test IPA
itself on a processor-sharing record. So this question may be small. It is
listed because ``untested'' and ``working'' are different claims, and
because a failure would be informative.

\subsection{What could go wrong}

Two distinct risks.

The mechanical risk: sliding a firing time in a multi-segment clock means
inverting the chain of segment laws. The inversion formulas exist and are
exercised by other models, but our processor-sharing law is a custom
distribution, and the inversion uses functions of it (the inverse survival
function in particular) that should be checked under differentiation.

The statistical risk: IPA is only valid when the statistic responds
smoothly to small shifts in the firing times. Time-integral statistics,
such as the average number in system, are smooth in this sense. But
processor sharing creates a subtlety: shifting one completion time changes
the sharing fraction over an interval, which changes \emph{other}
completion times through the re-declarations. The frozen-order replay does
propagate these effects through the segment chain, but whether the
propagation is complete for our construction is exactly what a test
against finite differences would reveal.

\subsection{What has to be decided}

Very little needs deciding in advance. The right move is an experiment:
run IPA on the processor-sharing model against finite differences, the
same comparison every other estimator passed. If it matches, we write down
the validity claim (which statistics, which couplings) and add the test to
the charter. If it does not match, the failure will point at either the
mechanical risk or the statistical risk, and that outcome feeds Question 2,
because both trace to the same underlying object: derivative flow through
re-declaration boundaries.

%==============================================================================
\section{Question 4: should the record store segment boundaries?}

\subsection{Where we are}

This question sits underneath the previous two.

Today the record stores only firings: clock, time, auxiliary values. It
does not store when clocks were enabled, and it does not store when clocks
were re-declared. All of that is \emph{reconstructed}: a bookkeeping walk
re-runs the model's own rules over the firing list and rebuilds each
clock's history, including its segments. Re-declarations are detected by
comparing distributions: after each firing, the walk rebuilds each active
clock's declared law and checks whether it changed.

\subsection{The case for reconstruction}

Reconstruction keeps the record minimal, and it forces the model's rules
to be the single source of truth. There is also a real audit benefit: when
the sampler's own notion of the clock history is compared against the
reconstruction and they agree, we have checked, on every run, that the
model rules and the simulation actually match. Storing the history instead
would remove that cross-check, because the record would just repeat what
the simulator believed.

\subsection{The case for storing}

Reconstruction has costs. It walks the whole record and rebuilds every
active clock's distribution after every firing, which is the most
expensive part of preparing a record for the estimators. It also rests on
an assumption: that a re-declaration always produces a visibly different
distribution. If a re-declaration ever produced an equal distribution ---
same law, arrived at for a different reason --- the walk would not see it.
Today that case cannot arise, because our only re-declaration source is a
change in the sharing fraction, which always changes the law. But the
assumption is load-bearing, and it is the kind of assumption that a future
feature could silently break.

Storing segment boundaries in the record would make preparation cheap and
the assumption unnecessary. The record would grow, and the audit would
weaken.

\subsection{What has to be decided}

This is a choice about where truth lives. Either the model's rules are
authoritative and the record is minimal evidence, or the record is
authoritative and the rules are one producer of it. The right answer may
depend on the outcomes of Questions 2 and 3: if re-declaration boundaries
become central to more estimators, the argument for recording them
directly gets stronger. The recommendation is to hold this question open
until at least the Question 3 experiment is done, and then decide with
that data in hand.

%==============================================================================
\section{Summary}

\begin{center}
\small
\begin{tabular}{@{}llll@{}}
\toprule
Question & Kind & First step & What it unlocks \\
\midrule
1. Auxiliary draws in clones & engineering & pick the stream keys & branching for SITA-style models \\
2. SPA with re-declaration & research & two-job worked example & SPA for processor sharing \\
3. IPA on PS records & experiment & finite-difference test & completes the estimator table \\
4. Segments in the record & architecture & wait for 2 and 3 & simpler, faster estimation prep \\
\bottomrule
\end{tabular}
\end{center}

A reasonable order: run the Question 3 experiment first, because it is
cheap and its outcome informs Questions 2 and 4. Then settle Question 1,
which is self-contained engineering. Then take up Question 2 with whatever
Question 3 revealed. Then revisit Question 4.

\end{document}
